\chapter{Spring Boot: Cỗ máy bên dưới mọi service}
\label{ch:springboot}

Bảy trong tám đơn vị triển khai của dự án này là ứng dụng Spring Boot.
Trước khi bất kỳ file YAML nào có thể có nghĩa, bạn cần ba ý tưởng: một
\emph{starter} kéo theo những gì, \emph{auto-configuration} quyết định
những gì thay cho bạn, và --- quan trọng nhất cho mọi thứ trong Phần~V và
VI --- chính xác \emph{Spring phân giải một thuộc tính cấu hình như thế
nào} khi cùng một key được định nghĩa ở nhiều nơi cùng lúc.

\section{Vấn đề Spring Boot giải quyết}

Spring cổ điển bắt bạn nối dây mọi thứ: một servlet container để triển
khai vào, cấu hình XML hoặc Java khai báo từng bean, và một danh sách dài
các dependency tương thích phiên bản với nhau. Spring Boot đảo ngược điều
đó. Một ứng dụng chỉ là một phương thức \code{main()} bình thường:

\begin{lstlisting}[language=Java]
@SpringBootApplication
public class CourseServiceApplication {
  public static void main(String[] args) {
    SpringApplication.run(CourseServiceApplication.class, args);
  }
}
\end{lstlisting}

\code{@SpringBootApplication} là ba annotation núp chung trong một chiếc
áo choàng: \code{@Configuration} (lớp này có thể định nghĩa bean),
\code{@ComponentScan} (tìm \code{@Service}, \code{@RestController},
\ldots{} trong package này và bên dưới), và \code{@EnableAutoConfiguration}
--- cái thú vị nhất.

\subsection{Starter}

Một starter như \code{spring-boot-starter-web} là một jar rỗng mà công
việc duy nhất là mang một danh sách dependency được tuyển chọn: Spring
MVC, Jackson, và một Tomcat \emph{nhúng}. ``Nhúng'' quan trọng về mặt vận
hành: server nằm bên trong jar, nên \code{java -jar course-service.jar}
là một web server hoàn chỉnh. Riêng sự thật đó là điều khiến các
container image ở Chương~13 đơn giản đến vậy --- không có Tomcat nào để
cài, chỉ cần cung cấp một JVM.

\subsection{Auto-configuration}

Lúc khởi động, Spring Boot duyệt qua một danh sách các cấu hình ứng viên
và áp dụng từng cái mà điều kiện của nó thỏa mãn. Các điều kiện đọc lên
kiểu như:

\begin{itemize}
  \item \emph{lớp này có trên classpath không?} (\code{@ConditionalOnClass}
  --- bạn đã thêm driver Postgres, nên một \code{DataSource} được cấu hình)
  \item \emph{người dùng đã tự định nghĩa bean này chưa?}
  (\code{@ConditionalOnMissingBean} --- bean tường minh của bạn luôn thắng)
  \item \emph{một thuộc tính có được đặt không?} (\code{@ConditionalOnProperty})
\end{itemize}

Auto-configuration là lý do \code{course-service} không chứa dòng code nào
tạo connection pool, Kafka producer, hay nối dây S3 client --- thêm
dependency cộng vài thuộc tính \emph{chính là} cấu hình. Khi hành vi
dường như xuất hiện từ hư không, hãy chạy ứng dụng với \code{--debug} một
lần: ``condition evaluation report'' liệt kê mọi auto-configuration được
áp dụng và bị bỏ qua, kèm lý do.

\section{Cấu hình ngoài: thứ tự phân giải thuộc tính}
\label{sec:property-order}

Đây là cơ chế chịu lực nhiều nhất trong cả cuốn sách. Mọi mẹo mà các
manifest Kubernetes ở Phần~V giở ra đều hoạt động nhờ nó.

Một thuộc tính Spring như \code{spring.datasource.url} có thể được định
nghĩa trong nhiều \emph{property source} (nguồn thuộc tính). Spring gộp
tất cả vào một \code{Environment}, và khi hai nguồn định nghĩa cùng một
key, nguồn cụ thể hơn thắng. Rút gọn về những nguồn dự án này thực sự
dùng, từ yếu nhất đến mạnh nhất:

\begin{enumerate}
  \item giá trị mặc định trong \code{application.yml} đóng gói trong jar
  \item cấu hình lấy từ config-server (Chương~2)
  \item biến môi trường của hệ điều hành
  \item tham số dòng lệnh
\end{enumerate}

\begin{decision}[vì sao cụm ghi đè bằng biến môi trường]
Cùng một image phải chạy không đổi trên laptop (Compose) lẫn trong k3s.
Nướng sự khác biệt giữa các môi trường vào jar nghĩa là mỗi môi trường
một image --- đúng thứ container sinh ra để tránh. Biến môi trường đứng
trên cả YAML đóng gói lẫn config-server, nên Deployment của Kubernetes có
thể trỏ lại một service sang \code{postgres-course:5432} mà không đụng
đến một file nào thuộc về lập trình viên. Một artifact, nhiều môi trường,
khác biệt được tiêm vào ở biên.
\end{decision}

\subsection{Ràng buộc nới lỏng (relaxed binding)}

Tên biến môi trường không thể chứa dấu chấm. Vì vậy Spring ánh xạ
\code{SPRING\_DATASOURCE\_URL} $\rightarrow$ \code{spring.datasource.url}:
viết hoa, dấu chấm thành gạch dưới. Đọc bất kỳ Deployment nào trong
\code{deploy/k8s/base/} và giờ bạn có thể giải mã mọi mục \code{env:}
thành ``cái này ghi đè key YAML kia.'' Ví dụ
\code{SPRING\_DATA\_MONGODB\_URI} trong \code{dictionary-service.yaml}
ghi đè \code{spring.data.mongodb.uri} từ config-server.

\subsection{Placeholder với giá trị mặc định}

Các file cấu hình của dự án dùng thêm một thành ngữ nữa:

\begin{lstlisting}[language=yaml]
app:
  jwt:
    secret: ${JWT_SECRET:mySecretKeyThatIs...}
\end{lstlisting}

\code{\$\{NAME:default\}} nghĩa là ``dùng biến môi trường nếu có, không
thì dùng giá trị mặc định.'' Đây là lý do một lập trình viên có thể chạy
auth-service mà không cần cài đặt gì, trong khi production tiêm secret
thật --- và là lý do ở Chương~16 Secret của Kubernetes chỉ cần cung cấp
\code{JWT\_SECRET} là nắm quyền kiểm soát giá trị này.

\section{File của dự án, có chú thích}

Mỗi service chỉ mang một \code{application.yml} tí hon; phần cấu hình thú
vị nằm trên config-server. Đây là file của course-service, đầy đủ:

\begin{lstlisting}[language=yaml]
spring:
  application:
    name: course-service            (*@\co{1}@*)
  config:
    import: optional:configserver:http://localhost:8888   (*@\co{2}@*)
eureka:
  client:
    service-url:
      defaultZone: http://localhost:8761/eureka/          (*@\co{3}@*)
\end{lstlisting}

\begin{description}
  \item[\co{1}] Danh tính của service. Config-server dùng nó để chọn file
  cần phục vụ (\code{configurations/course-service.yml}); Eureka dùng nó
  làm tên đăng ký; gateway định tuyến tới nó bằng
  \code{lb://course-service}. Một chuỗi ký tự buộc cả ba hệ thống lại
  với nhau.
  \item[\co{2}] Import cấu hình từ xa lúc bootstrap.
  \code{optional:} nghĩa là ``cứ khởi động dù config-server có sập'' ---
  dễ chịu khi dev, và đáng bàn ở production (một service khởi động với
  nửa cấu hình có thể hỏng theo những cách lạ lùng hơn một service từ
  chối khởi động). Giá trị mặc định \code{localhost:8888} được ghi đè bởi
  \code{SPRING\_CONFIG\_IMPORT} trong cả Compose lẫn Kubernetes.
  \item[\co{3}] Eureka ở đâu --- lại một giá trị mặc định localhost, lại
  được ghi đè theo từng môi trường bằng
  \code{EUREKA\_CLIENT\_SERVICEURL\_DEFAULTZONE}.
\end{description}

\section{Profile}

Một \emph{profile} Spring là một biến thể cấu hình có tên.
\code{spring.profiles.active=native} trên config-server (Chương~2) là
chỗ dùng profile duy nhất của dự án cho đến giờ, nhưng cơ chế này tổng
quát hơn: \code{application-prod.yml} được nạp đè lên
\code{application.yml} khi profile \code{prod} đang kích hoạt, cho phép
một jar mang nhiều tính cách. Kubernetes làm giảm nhu cầu đó --- thay vào
đó môi trường tiêm sự khác biệt vào --- đó là lý do dự án này dựa vào
biến môi trường, không phải file profile, cho thay đổi theo môi trường.

\section{Actuator}

\code{spring-boot-starter-actuator} thêm các endpoint vận hành dưới
\code{/actuator}: \code{/health} (readiness tổng hợp của DB, đĩa, các
kiểm tra tùy chỉnh), \code{/info}, \code{/metrics}, \code{/prometheus}.
Hai sự thật quan trọng lúc này:

\begin{itemize}
  \item Starter phải có trên classpath. Cấu hình exposure mà không có
  \code{spring-boot-starter-actuator} thì vô tác dụng --- dự án đã mắc
  đúng lỗi đó ở bốn service và phát hiện ra qua một probe Kubernetes thất
  bại (Chương~17).
  \item Endpoint phải được \emph{expose}
  (\code{management.endpoints.web.exposure.include}) --- dự án expose
  \code{health,info,metrics,prometheus}.
  \item Expose không phải là \emph{phân quyền}: nếu Spring Security bảo vệ
  \code{/actuator/**}, endpoint tồn tại nhưng trả lời 401. Chính tương
  tác này quyết định service nào dùng được probe Kubernetes nào ở
  Chương~17 --- dictionary-service cho phép actuator và được probe qua
  \code{/actuator/health}; auth-service thì không, và phải được probe qua
  \code{/v3/api-docs} thay thế.
\end{itemize}

\section{Bên ngoài dự án}

Hai ví dụ tương phản của cùng một bộ máy thuộc tính:

\emph{Một ứng dụng twelve-factor kiểu Heroku} cấu hình mọi thứ bằng biến
môi trường --- \code{DATABASE\_URL}, \code{PORT} --- và mang một
\code{application.yml} chỉ toàn placeholder \code{\$\{...\}}. Thứ tự phân
giải của Spring khiến phong cách đó trở thành bản địa.

\emph{Một ngân hàng chạy một artifact qua năm môi trường}
(dev/test/uat/preprod/prod) thường ghép một config-server dựa trên git
repository với profile cho từng môi trường:
\code{course-service-uat.yml} nằm trong một repo được kiểm toán, và việc
thăng cấp là một lần merge git. Config-server classpath native của dự án
(Chương~2) là phiên bản một-lập-trình-viên của cùng hình dạng đó.

\begin{pitfall}[giá trị ghi đè của tôi bị bỏ qua]
Triệu chứng: bạn đặt một biến môi trường nhưng service vẫn dùng giá trị
cũ. Các nguyên nhân thường gặp, theo thứ tự tần suất: tên biến không ánh
xạ tới key bạn nghĩ (kiểm tra cách viết theo relaxed binding ---
\code{SERVICEURL} trong key Eureka là một từ, một cái tên bất quy tắc
phải chép chính xác); giá trị cũng được đặt bởi một nguồn \emph{mạnh
hơn} (một tham số dòng lệnh); hoặc service đọc key lúc build thay vì từ
Environment. Chẩn đoán bằng endpoint actuator \code{/actuator/env} (khi
được expose), nó liệt kê mọi property source và nguồn nào đã thắng.
\end{pitfall}

\begin{quiz}
\begin{enumerate}
  \item Một key được định nghĩa trong \code{application.yml} của jar,
  trong file của config-server cho service đó, và dưới dạng biến môi
  trường trong Deployment. Giá trị nào thắng, và vì sao thứ tự đó là đúng
  cho container?
  \item \code{@SpringBootApplication} gộp ba annotation nào, và annotation
  nào trong số đó giải thích vì sao xóa một dependency không dùng đến có
  thể làm thay đổi hành vi lúc chạy?
  \item \code{\$\{GOOGLE\_CLIENT\_ID:placeholder\}} --- biểu thức này phân
  giải thành gì khi biến chưa được đặt, và điều đó cho auth-service tính
  chất vận hành nào lúc khởi động?
  \item Vì sao expose một endpoint actuator không đảm bảo rằng một probe
  Kubernetes dùng được nó?
\end{enumerate}
\end{quiz}
